\subsubsection{Deux plaquettes animées : protocole et accès aux fichiers}
Les fichiers GIF accompagnent le PDF dans le sous-dossier \texttt{animations}. Le PDF présente des images fixes et des liens vers les animations externes ; il ne joue pas lui-même les GIF. Les liens supposent de conserver ce sous-dossier à côté du rapport.
\begin{itemize}
\item \href{animations/plaquette_cycle_thermique.gif}{Ouvrir la plaquette du cycle chauffage--refroidissement}.
\item \href{animations/plaquette_relaxation_isotherme.gif}{Ouvrir la plaquette de relaxation à température fixe}.
\end{itemize}
Chaque plaquette contient neuf couches atomiques exactes : trois couches parallèles pour chacune des orientations $(001)$, $(110)$ et $(111)$. Les positions des coupes et le cadrage restent fixes ; Au est en or et Cu en bleu. Pour $(001)$, on suit $z/a=2,2.5,3$ ; pour $(110)$, $(x+y)/a=5,5.5,6$ ; pour $(111)$, $(x+y+z)/a=7,8,9$. Ces trois orientations sont des vues complémentaires du même volume, pas trois simulations distinctes.
Le GIF thermique comporte 93 images : l'état initial puis les 92 fins de palier du premier cycle, de 300 à 4000 K puis en sens inverse. Les événements à l'intérieur d'un palier ne sont pas sauvegardés dans cette campagne historique. Le compteur cumule équilibration et production.
Le GIF isotherme comporte 301 états. Il part de l'état final à 4000 K du chauffage du cycle 1, puis impose immédiatement 1000 K. Un calcul dédié de 6000 balayages, graine 20261104, sauvegarde les occupations tous les 20 balayages. Aucun transitoire n'est éliminé. La première image est l'état chaud juste après le changement de température, avant toute relaxation ; l'étiquette 1000 K indique la consigne. Les 200 premiers balayages sont ralentis à la lecture, sans modifier les valeurs du compteur.
\subsubsection{Ce que montrent les animations et ce qu’elles ne démontrent pas}
Le cycle thermique rend visible le passage entre un motif ordonné et une occupation désordonnée. La variante retrouvée au refroidissement peut différer de celle du départ : un changement de variante n'est pas, à lui seul, une perte d'ordre. Une seule couche peut paraître pauvre en Au parce que le sous-réseau privilégié se trouve dans une couche voisine.
Dans la relaxation, $q$ passe de $0.049$ à $0.944$ entre 0 et 6000 balayages. Les premiers états enregistrés avec $q\geq0.8$ et $q\geq0.9$ apparaissent respectivement à 80 et 100 balayages (résolution de capture : 20 balayages). L'évolution ultérieure fluctue et n'est pas strictement monotone.
Cette trajectoire montre une mise en ordre rapide de la petite boîte. Elle ne suffit pas à identifier un rayon critique, un temps d'attente de nucléation ou une loi de croissance de domaines. Il faudrait définir un domaine local, suivre ses interfaces et répéter les trempes avec plusieurs graines et tailles. Les échanges Au--Cu sont globaux : les balayages ne constituent pas une durée physique ni une cinétique de diffusion atomique calibrée.
\fig{animations_01_evolution}{1}{Ordre global des configurations présentes dans les deux GIF et agrandissement du début de la relaxation.}
Les plaquettes fixes initiale et finale sont disponibles avec les GIF dans \texttt{results/thermal\_cycle\_L6/visualisation/animations} (fichiers \texttt{plaquette\_relaxation\_isotherme\_debut.pdf} et \texttt{plaquette\_relaxation\_isotherme\_fin.pdf}). Elles ne sont pas reproduites une seconde fois comme grandes figures de neuf panneaux.
